Compte tenu de la variation d'une année à l'autre observée de 2015 à 2023, on ne peut attendre de ce dossier qu'il détecte un changement lié à la piste dans les dons de la paroisse inférieur à environ 15 à 26~\% des dons d'une année contre l'histoire de la paroisse elle-même, ou à environ 14 à 16 points contre la variation médiane des douze paroisses voisines (deux écarts-types des variations d'avant la piste, le chiffre le plus bas sans les trois transitions COVID~; section~\ref{sec:threshold}). Comme la piste a couvert sept des douze mois de 2024, un effet de la piste aurait dû retrancher environ 25 à 45~\% des dons annuels de la paroisse, et environ 21 à 50~\% de ses revenus totaux, pour se détacher dans la déclaration de 2024~; les variations observées en 2024 ($-$6,1~\% sur les dons, $-$6,9 points contre les voisines) sont bien en deçà de ces seuils. La déclaration de 2025, une année complète, abaisserait le seuil à environ 14 à 26~\% des dons, et deux années complètes à environ 10 à 18~\%.
